\documentclass[12pt]{article}
\usepackage[utf8]{inputenc}
\usepackage[T1]{fontenc}
\usepackage{hyperref}
\title{Filipino Sign Language Language Resources for K-12 Curriculum Development}
\author{Compiled by Claude Sonnet 5.5 (claude-sonnet-5-5) for Andr\'as Kornai}
\date{}
\begin{document}
\maketitle

\section{Introduction}
Filipino Sign Language (FSL; ISO 639-3 psp) is the national sign language of the Philippines under Republic Act 11106. This dossier lists what turned up in a brief web search. Several items were seen only as search-result titles and summaries, so check each before use.

\section{Primary Educational Resources}
\begin{itemize}
\item DepEd, Kindergarten and Grades 1 to 3 Mother Tongue Most Essential Learning Competencies for Filipino Sign Language: \url{https://lrmds.deped.gov.ph/detail/18648}
\item Benilde SDEAS FSL courses (graded learning programme): \url{https://sdeas.benilde.edu.ph/courses/}
\end{itemize}

\section{Linguistic and Reference Materials}
\begin{itemize}
\item Benilde, field documentation of FSL vocabulary and phrases; the search summary links it to the FSL Student Handbook Level 1, with grammar and a dictionary of basic communication: \url{https://sdgs.benilde.edu.ph/field-documentation-of-filipino-sign-language-vocabulary-and-phrases/}
\item Benilde, five facts about FSL: \url{https://sdeas.benilde.edu.ph/fivefactsaboutfsl/}
\item Benilde SDGs page: \url{https://sdgs.benilde.edu.ph/?p=8965}
\item UP College of Education, a COVID-19 dictionary for Filipino children with FSL videos (as described in the search summary): \url{https://educ.upd.edu.ph/?p=2719}
\end{itemize}

\section{Religious and Community Resources}
\begin{itemize}
\item Philippine Federation of the Deaf, Filipino Sign Language Act page: \url{https://pfdeaf.org/filipino-sign-language-act/}
\item Philippine Federation of the Deaf, education document: \url{https://pfdeaf.org/wp-content/uploads/2023/03/education.pdf}
\end{itemize}
No scripture resources were found.

\section{Digital and Audio Resources}
Only the FSL video links in the UP dictionary and Benilde materials were found. No dedicated corpus or software was identified.

\section{Curriculum Development Context}
RA 11106 declares FSL the national sign language and the language of instruction in deaf education, and requires it as a separate subject for deaf learners. The text of the Act was listed at \url{https://chr-observatories.uwazi.io/en/entity/8ydk0adx056}. DepEd competencies were found for K--3 only. No complete national FSL dictionary or full K-12 curriculum was found.
\end{document}